% ===========================================================================
% Poste 2 — RQAP
% ===========================================================================
\poste{2}{Régime québécois d'assurance parentale (RQAP)}{QC\_rqap}

\pdesc Cotisation sur le revenu de travail finançant les prestations de
maternité, de paternité, parentales et d'adoption du régime québécois.

\pobj Soutenir financièrement les parents à la naissance ou à l'adoption
d'un enfant, en remplaçant une partie de leur revenu pendant le congé.

\pregle Contrairement au RRQ, le RQAP ne comporte \emph{aucune exemption} :
dès que le revenu assurable dépasse le seuil minimal $s$, la cotisation
s'applique au \emph{plein} revenu, plafonné au maximum assurable $M$ :
\[
\text{cotisation} =
\begin{cases}
0 & \text{si } r \le s \\
\tau \times \min(r,\, M) & \text{sinon}
\end{cases}
\]
où $\tau$ est le taux de cotisation (part de l'employé). La cotisation du
ménage est la somme des cotisations des adultes ; les ménages retraités ne
cotisent pas. À l'impôt fédéral, la cotisation RQAP donne droit à un crédit
non remboursable (poste~19).

\emph{Point de contrôle.} Cotisation maximale 2025 :
$0{,}00494 \times \D{98000} = \D{484.12}$ ; 2026 :
$0{,}0043 \times \D{103000} = \D{442.90}$.

\pparams

\begin{center}
\small
\begin{tabular}{l r r}
\toprule
Paramètre & 2025 & 2026 \\
\midrule
Revenu maximal assurable ($M$)      & \D{98000}   & \D{103000} \\
Revenu assurable minimal ($s$)      & \D{2000}    & \D{2000} \\
Taux de cotisation ($\tau$)         & \pc{0.494}  & \pc{0.430} \\
\bottomrule
\end{tabular}
\end{center}

La baisse du taux en 2026 (de \pc{0.494} à \pc{0.430}) est le second
allègement de cotisation visible dans l'écart 2025--2026.

% ============================================================================
% GÉNÉRÉ par scripts/exemples-guide.ts — NE PAS ÉDITER À LA MAIN.
% Régénérer : npm run exemples (chaque chiffre est vérifié contre le moteur).
% ============================================================================
\pexemples Taux unique de \pc{0.494} sur le revenu assurable, \emph{sans
exemption} : dès que le revenu dépasse \D{2000}, la cotisation porte sur
le plein montant, plafonné au maximum assurable de \D{98000}.

\medskip\noindent\textbf{M2 --- personne vivant seule, 30~ans, revenu de travail de \D{9000}.}
\begin{align*}
\text{cotisation} &= \num{9000} \times \pc{0.494} = \num{44.46}
\end{align*}
Le revenu dépasse le seuil de \D{2000} : la cotisation s'applique dès le
premier dollar, soit \D{44.46}.

\medskip\noindent\textbf{M4 --- personne vivant seule, 40~ans, revenu de travail de \D{50000}.}
\begin{align*}
\text{cotisation} &= \num{50000} \times \pc{0.494} = \num{247}
\end{align*}
Cotisation RQAP : \D{247}.

\medskip\noindent\textbf{M5 --- personne vivant seule, 45~ans, revenu de travail de \D{100000}.}
Le revenu (\D{100000}) excède le maximum assurable : la cotisation
est plafonnée.
\begin{align*}
\text{cotisation} &= \min(\num{100000},\ \num{98000}) \times \pc{0.494} = \num{484.12}
\end{align*}
Cotisation RQAP maximale 2025 : \D{484.12}.


\prefs Loi sur l'assurance parentale (RLRQ, c.~A-29.011) ; Conseil de
gestion de l'assurance parentale ; Revenu Québec.
